\chapter{State of the Art}

\section{Computer Vision for Queue Monitoring}
Object detection and tracking in dynamic scenes are fundamental building blocks for intelligent monitoring.
Modern approaches mainly rely on convolutional-neural-network detectors and multi-object tracking algorithms.

\section{People Detection}
The YOLO family offers a strong trade-off between accuracy and latency for real-time applications.
In this project, YOLO26 variants are selected based on operational needs (speed vs. precision).

\section{Multi-Object Tracking}
Tracking keeps a consistent identity for each person across consecutive frames.
ByteTrack is chosen for its robustness and straightforward integration into a real-time pipeline.

\section{Queueing Modeling}
Queueing theory provides a mathematical framework to estimate service-system performance.
The M/M/1 model is selected as a first approximation, under the stability condition $\lambda < \mu$.

\section{Uncertainty Quantification}
Point estimates alone can be misleading in noisy environments.
Bayesian approaches, especially Gamma distributions for Poisson rates, make it possible to produce interpretable credible intervals.

\section{Critical Synthesis}
Existing work often focuses on detection or counting, with less emphasis on:
\begin{itemize}
    \item explicit waiting-time estimation,
    \item uncertainty quantification,
    \item operational integration (alerts, dashboard, webhook).
\end{itemize}
This final year project is positioned precisely at this intersection.
